% ============================================================
%  Presentación Beamer "UNI Oscuro" — Conversión de Energía (Ciclos)
%  Termodinámica MN121 · compilar con xelatex (dos pasadas)
% ============================================================
\documentclass[aspectratio=169,10pt]{beamer}

\usepackage{fontspec}
\usepackage[spanish,es-noshorthands]{babel}
\usepackage{tikz}
\usetikzlibrary{shapes.geometric,positioning}
\usepackage[most]{tcolorbox}
\usepackage{fontawesome5}
\usepackage{booktabs,colortbl,array,ragged2e,graphicx}
\usepackage{mathtools}
\graphicspath{{img/}}

% ---------- Paleta ----------
\definecolor{FondoOscuro}{HTML}{040812}
\definecolor{AzulPanel}{HTML}{0C111C}
\definecolor{Granate}{HTML}{660F1A}
\definecolor{GranateOscuro}{HTML}{3D0710}
\definecolor{GranateVelo}{HTML}{381520}
\definecolor{Teal}{HTML}{00A6A6}
\definecolor{TealOscuro}{HTML}{01565B}
\definecolor{Dorado}{HTML}{D4AF37}
\definecolor{Naranja}{HTML}{B46A35}
\definecolor{Cian}{HTML}{00F2FF}
\definecolor{Crema}{HTML}{FAF9F7}
\definecolor{CremaGris}{HTML}{F6F5F3}
\definecolor{TintaTexto}{HTML}{2F3E46}
\definecolor{TealSuave}{HTML}{F2FBFB}
\definecolor{GrisLinea}{HTML}{E2E1E0}
\definecolor{IconoTenue}{HTML}{0B2430}

% ---------- Fuentes ----------
\setsansfont[
  BoldFont=FiraSans-SemiBold.otf,
  ItalicFont=FiraSans-Italic.otf,
  BoldItalicFont=FiraSans-SemiBoldItalic.otf
]{FiraSans-Regular.otf}
\setmonofont[Scale=0.9]{FiraCode-Regular.ttf}
\usefonttheme{professionalfonts}
\renewcommand{\familydefault}{\sfdefault}

% ---------- METADATOS ----------
\newcommand{\sellocorto}{UNI | FIM | MN121}
\newcommand{\pietitulo}{Conversión de Energía · Ciclos}

% ---------- Ajustes globales ----------
\setbeamertemplate{navigation symbols}{}
\setbeamersize{text margin left=8mm, text margin right=8mm}
\setbeamercolor{normal text}{fg=TintaTexto, bg=Crema}
\setbeamercolor{structure}{fg=Granate}
\setbeamercolor{alerted text}{fg=Granate}
\setbeamercolor{example text}{fg=TealOscuro}
\setbeamertemplate{itemize item}{\color{Teal}\raisebox{0.4ex}{\scriptsize\faCaretRight}}
\setbeamertemplate{itemize subitem}{\color{Granate}\raisebox{0.4ex}{\tiny\faCaretRight}}
\setbeamerfont{frametitle}{series=\bfseries,size=\large}

% ---------- Fondo ----------
\setbeamertemplate{background canvas}{%
\begin{tikzpicture}
  \fill[Crema] (0,0) rectangle (16,9);
  \draw[white,       line width=1.3pt] (0,2.1) .. controls (4.2,3.5) and (9.5,1.1) .. (16,3.1);
  \draw[GrisLinea!60,line width=0.9pt] (0,1.0) .. controls (5.0,2.3) and (10.5,0.3) .. (16,1.7);
  \draw[white,       line width=1.3pt] (0,5.3) .. controls (5.2,6.8) and (11.0,4.5) .. (16,6.3);
  \draw[GrisLinea!60,line width=0.9pt] (0,7.3) .. controls (6.0,8.5) and (11.5,6.5) .. (16,7.9);
\end{tikzpicture}}

% ---------- Cabecera ----------
\setbeamertemplate{frametitle}{%
\begin{tikzpicture}[remember picture, overlay]
  \fill[FondoOscuro] (current page.north west) rectangle ([yshift=-0.85cm]current page.north east);
  \fill[Granate] (current page.north west)
      -- ([xshift=7.6cm]current page.north west)
      -- ([xshift=6.7cm,yshift=-0.85cm]current page.north west)
      -- ([yshift=-0.85cm]current page.north west) -- cycle;
  \fill[Dorado] ([yshift=-0.91cm]current page.north west)
      rectangle ([xshift=6.7cm,yshift=-0.85cm]current page.north west);
  \fill[Teal] ([xshift=6.7cm,yshift=-0.97cm]current page.north west)
      rectangle ([yshift=-0.85cm]current page.north east);
  \fill[Teal] ([xshift=-0.42cm]current page.north east)
      rectangle ([xshift=-0.28cm,yshift=-0.30cm]current page.north east);
  \node[anchor=west, text=white] at ([xshift=0.55cm,yshift=-0.43cm]current page.north west)
      {\large\bfseries\insertframetitle};
\end{tikzpicture}%
\vspace*{0.42cm}}

% ---------- Pie ----------
\setbeamertemplate{footline}{%
\begin{tikzpicture}[remember picture, overlay]
  \fill[CremaGris] ([yshift=0.5cm]current page.south west) rectangle (current page.south east);
  \draw[Teal, line width=0.7pt] ([yshift=0.5cm]current page.south west) -- ([yshift=0.5cm]current page.south east);
  \fill[Granate] ([xshift=-0.34cm]current page.south east) rectangle ([yshift=0.5cm]current page.south east);
  \node[anchor=west, text=FondoOscuro] at ([xshift=0.55cm,yshift=0.25cm]current page.south west)
      {\scriptsize \faUniversity\hspace{1mm}\textbf{\sellocorto}\hspace{1.4mm}{\color{Teal}\faAngleRight}\hspace{1.4mm}\textit{\pietitulo}};
  \node[anchor=east, text=FondoOscuro] at ([xshift=-0.55cm,yshift=0.25cm]current page.south east)
      {\scriptsize\textbf{\insertframenumber}\hspace{0.8mm}{\tiny\inserttotalframenumber}};
\end{tikzpicture}}

% ---------- Tarjetas ----------
\newtcolorbox{cardidea}[3][]{enhanced, colback=white, frame hidden, boxrule=0pt,
  arc=1.6mm, left=3mm, right=3mm, top=2.2mm, bottom=2.4mm,
  borderline west={2.6pt}{0pt}{Granate},
  drop fuzzy shadow={black!30},
  fontupper=\small, coltext=TintaTexto,
  before upper={{\color{Granate}#3}\hspace{1.6mm}{\normalsize\bfseries\color{FondoOscuro}#2}\par\vspace{1.3mm}}}

\newtcolorbox{cardgear}[3][]{enhanced, colback=TealSuave, frame hidden, boxrule=0pt,
  arc=1.6mm, left=3mm, right=3mm, top=2.2mm, bottom=2.4mm,
  borderline west={2.6pt}{0pt}{Teal},
  drop fuzzy shadow={black!30},
  fontupper=\small, coltext=TintaTexto,
  before upper={{\color{TealOscuro}#3}\hspace{1.6mm}{\normalsize\bfseries\color{FondoOscuro}#2}\par\vspace{1.3mm}}}

\newtcolorbox{cardnota}[1][\faExclamationTriangle]{enhanced, colback=white,
  colframe=GrisLinea, boxrule=0.5pt, arc=1.4mm,
  left=3mm, right=3mm, top=1.8mm, bottom=2mm,
  drop fuzzy shadow={black!25},
  fontupper=\small, coltext=TintaTexto,
  before upper={{\color{FondoOscuro}#1}\hspace{1.6mm}}}

\newtcolorbox{cardlista}{enhanced, colback=white, frame hidden, boxrule=0pt,
  arc=1.2mm, left=3.5mm, right=2.5mm, top=2.4mm, bottom=2.4mm,
  borderline west={3pt}{0pt}{FondoOscuro},
  drop fuzzy shadow={black!30},
  fontupper=\small, coltext=Granate}

% ---------- Tabla de etapas ----------
\newcommand{\etapastitulo}{\par{\centering{\color{Granate}\faTable}\hspace{1.6mm}{\bfseries\color{FondoOscuro}Etapas del ciclo}\par}\vspace{0.2mm}}
\newenvironment{etapas}{%
  \begingroup\scriptsize\renewcommand{\arraystretch}{1.12}\arrayrulecolor{Granate}
  \begin{center}
  \begin{tabular}{@{}>{\centering\arraybackslash\bfseries}p{1.2cm} >{\RaggedRight}p{4.5cm} >{\RaggedRight\arraybackslash}p{6.6cm}@{}}
  \toprule
  {\color{FondoOscuro}Etapa} & {\color{FondoOscuro}\bfseries Proceso} & {\color{FondoOscuro}\bfseries Ecuación clave}\\
  \midrule}%
 {\bottomrule
  \end{tabular}
  \end{center}\endgroup}

% ---------- Leyendas ----------
\newcommand{\tcap}[1]{\begin{center}\vspace{-1mm}{\scriptsize{\color{Granate}\bfseries Tabla.} {\color{TintaTexto!75}#1}}\end{center}\vspace{-2.5mm}}
\newcommand{\fcap}[1]{{\par\centering\scriptsize{\color{Granate}\bfseries Figura.} {\color{TintaTexto!75}#1}\par}}
\newcommand{\fsub}[1]{{\par\centering\scriptsize\color{TintaTexto!75}#1\par}}

% ---------- Portada de sección ----------
\newcommand{\portadaseccion}[4]{%
\begin{frame}[plain]
\begin{tikzpicture}[remember picture, overlay]
  \fill[FondoOscuro] (current page.south west) rectangle (current page.north east);
  \node[color=IconoTenue] at ([xshift=-3.4cm,yshift=0.3cm]current page.east)
      {\fontsize{62}{62}\selectfont #4};
  \fill[GranateVelo] (current page.north west)
      -- ([xshift=8.6cm]current page.north west)
      -- ([xshift=11.3cm]current page.south west)
      -- (current page.south west) -- cycle;
  \fill[GranateOscuro, opacity=0.75] (current page.north west)
      -- ([xshift=6.4cm]current page.north west)
      -- ([xshift=9.1cm]current page.south west)
      -- (current page.south west) -- cycle;
  \draw[Teal, line width=1.5pt] ([xshift=8.6cm]current page.north west)
      -- ([xshift=11.3cm]current page.south west);
  \draw[Naranja, line width=0.9pt] ([xshift=7.4cm]current page.north west)
      -- ([xshift=10.1cm]current page.south west);
  \node[anchor=west, text=white] at ([xshift=1.1cm,yshift=0.75cm]current page.west)
      {\fontsize{24}{28}\selectfont\bfseries #1.~#2};
  \draw[white, line width=0.9pt] ([xshift=1.15cm,yshift=0.12cm]current page.west) -- ++(4.8cm,0);
  \node[anchor=west, text=white] at ([xshift=1.15cm,yshift=-0.55cm]current page.west)
      {{\color{white}\faAngleDoubleRight}\hspace{1.5mm}\small #3};
\end{tikzpicture}
\end{frame}}

\begin{document}

% ============================================================
%  PORTADA
% ============================================================
\begin{frame}[plain]
\begin{tikzpicture}[remember picture, overlay]
  \fill[FondoOscuro] (current page.south west) rectangle (current page.north east);
  \node[color=AzulPanel!80!white] at ([xshift=-4.0cm,yshift=-1.2cm]current page.north east)
      {\fontsize{22}{22}\selectfont $\eta_t=1-\dfrac{q_{sal}}{q_{ent}}$};
  \node[color=AzulPanel!80!white] at ([xshift=3.2cm,yshift=1.1cm]current page.south west)
      {\fontsize{20}{20}\selectfont $w_{neto}=q_{ent}-q_{sal}$};
  \node[color=AzulPanel!80!white] at ([xshift=-2.6cm,yshift=2.6cm]current page.south east)
      {\fontsize{18}{18}\selectfont $\mathrm{COP}=\dfrac{q_L}{w}$};
  \draw[Dorado, line width=0.8pt] ([xshift=-2.7cm,yshift=-3.4cm]current page.north east)
      -- ++(-1.5,-1.5) -- ++(-1.8,0);
  \fill[Cian] ([xshift=-6.0cm,yshift=-4.9cm]current page.north east) circle (0.055cm);
  \draw[Teal, line width=0.8pt] ([xshift=-1.6cm,yshift=-5.6cm]current page.north east)
      -- ++(-1.2,-1.2) -- ++(-2.2,0);
  \fill[Cian] ([xshift=-5.0cm,yshift=-6.8cm]current page.north east) circle (0.055cm);
  \fill[Cian] ([xshift=-1.6cm,yshift=-5.6cm]current page.north east) circle (0.05cm);
  \node[regular polygon, regular polygon sides=6, minimum size=2.5cm,
        draw=Teal, line width=1.3pt, shape border rotate=30]
        at ([xshift=-3.1cm,yshift=-0.2cm]current page.east) {};
  \node[color=Teal] at ([xshift=-3.1cm,yshift=-0.2cm]current page.east)
      {\fontsize{26}{26}\selectfont \faIndustry};
  \node[anchor=north west, text=white] at ([xshift=1.0cm,yshift=-0.85cm]current page.north west)
      {\footnotesize\bfseries UNIVERSIDAD NACIONAL DE INGENIERIA\ \textbar\ FIM\ \textbar\ MN121};
  \node[anchor=north west, text=white, align=left,
        font=\fontsize{19.5}{24}\selectfont\bfseries] at ([xshift=0.95cm,yshift=-1.55cm]current page.north west)
      {Conversión de Energía:\\ Ciclos Termodinámicos};
  \node[anchor=north west, text=Teal] at ([xshift=1.0cm,yshift=-4.05cm]current page.north west)
      {\normalsize\itshape Vapor\ \textperiodcentered\ Gas\ \textperiodcentered\ Combustión\ \textperiodcentered\ Refrigeración};
  \draw[Dorado, line width=0.6pt] ([xshift=1.0cm,yshift=2.45cm]current page.south west) -- ++(6.2cm,0);
  \node[anchor=south west, text=white] at ([xshift=1.0cm,yshift=1.75cm]current page.south west)
      {\normalsize\bfseries Termodinámica --- Ciclos de Potencia y Refrigeración};
  \node[anchor=south west, text=white] at ([xshift=1.0cm,yshift=1.25cm]current page.south west)
      {\small Curso: MN121\ \textbar\ Termodinámica --- FIM};
  \node[anchor=south west, text=white!70!FondoOscuro] at ([xshift=1.0cm,yshift=0.78cm]current page.south west)
      {\footnotesize Rankine \textperiodcentered\ Brayton \textperiodcentered\ Otto \textperiodcentered\ Diesel \textperiodcentered\ VCR \textperiodcentered\ Lima, 2026};
\end{tikzpicture}
\end{frame}

% ============================================================
%  RESUMEN
% ============================================================
\begin{frame}{Resumen}
\begin{cardidea}{Síntesis}{\faLightbulb}
Un \textbf{ciclo termodinámico} convierte calor en trabajo (potencia) o mueve calor consumiendo trabajo (refrigeración). Todo ciclo de potencia obedece $\eta_t = 1-\frac{q_{sal}}{q_{ent}}$ y $w_{neto}=q_{ent}-q_{sal}$; los de refrigeración se miden con el $\mathrm{COP}$. De \textbf{cada ciclo} damos su \textbf{definición}, sus \textbf{etapas} con la ecuación clave de cada transición y sus \textbf{ecuaciones generales}, más un diagrama y un ejemplo numérico.
\end{cardidea}
\vspace{2mm}
\begin{columns}[T]
\begin{column}{0.485\textwidth}
\begin{cardgear}{Ciclos de potencia}{\faFire}
Producen $w_{neto}>0$. Se mejoran subiendo la temperatura media de adición de calor: \textbf{recalentamiento}, \textbf{regeneración}, mayor $r_P$ o $r$.
\end{cardgear}
\end{column}
\begin{column}{0.485\textwidth}
\begin{cardgear}{Ciclos inversos}{\faSnowflake}
Consumen $w_C$ para bombear calor de frío a caliente. Métrica: $\mathrm{COP}_R=q_L/w_C$ (refrigerador) y $\mathrm{COP}_{HP}=\mathrm{COP}_R+1$ (bomba de calor).
\end{cardgear}
\end{column}
\end{columns}
\end{frame}

% ============================================================
%  CLASIFICACIÓN
% ============================================================
\begin{frame}{Clasificación de los ciclos}
\begin{columns}[T]
\begin{column}{0.5\textwidth}
\centering
\includegraphics[width=\linewidth,height=4.6cm,keepaspectratio]{ciclos_clasificacion}\\[0.5mm]
\fsub{Mapa de familias de ciclos termodinámicos.}
\end{column}
\begin{column}{0.46\textwidth}
\begin{cardidea}{Ejes de clasificación}{\faInfoCircle}
\textbf{Objetivo:} potencia vs.\ refrigeración.\\[0.6mm]
\textbf{Fluido:} condensable (vapor) vs.\ gas.\\[0.6mm]
\textbf{Combustión:} externa (Rankine/Brayton cerrado) vs.\ interna (Otto/Diesel).
\end{cardidea}
\vspace{1.5mm}
\begin{cardlista}
1. Ciclos de Vapor --- Rankine\\[0.6mm]
2. Ciclos de Gas --- Brayton\\[0.6mm]
3. Combustión Interna --- Otto / Diesel\\[0.6mm]
4. Refrigeración --- Compresión de Vapor
\end{cardlista}
\end{column}
\end{columns}
\end{frame}

% ============================================================
%  SECCIÓN 1 — VAPOR
% ============================================================
\portadaseccion{1}{Ciclos de Vapor --- Rankine}{caldera · turbina · condensador · bomba}{\faFire}

% ---- Rankine Simple: definición y etapas ----
\begin{frame}{Rankine Simple --- definición y etapas}
\begin{cardidea}{Definición}{\faInfoCircle}
Versión \textbf{ideal de cuatro estados} del ciclo de vapor: turbina y bomba isentrópicas, caldera y condensador isobáricos. Es el ciclo base de las \textbf{centrales termoeléctricas}; el fluido (agua) cambia de fase. Las irreversibilidades reales se modelan con $\eta_T,\eta_P$.
\end{cardidea}
\vspace{1.4mm}
\etapastitulo
\begin{etapas}
$1\to2$ & Bombeo isentrópico ($P_L\to P_H$)      & $w_B=h_2-h_1=v_1(P_H-P_L)$\\
$2\to3$ & Adición de calor isobárica (caldera)    & $q_H=h_3-h_2$\\
$3\to4$ & Expansión isentrópica (turbina)         & $w_T=h_3-h_4$\\
$4\to1$ & Rechazo de calor isobárico (condensador)& $q_L=h_4-h_1$\\
\end{etapas}
\end{frame}

% ---- Rankine Simple: ecuaciones y ejemplo ----
\begin{frame}{Rankine Simple --- ecuaciones y ejemplo}
\begin{columns}[T]
\begin{column}{0.46\textwidth}
\centering
\includegraphics[width=\linewidth,height=4.2cm,keepaspectratio]{rankine_esquema_planta}\\[0.5mm]
\fsub{Esquema de la planta: caldera, turbina, condensador y bomba.}
\end{column}
\begin{column}{0.5\textwidth}
\begin{cardidea}{Ecuaciones generales}{\faSuperscript}
$\eta_{th}=\dfrac{(h_3-h_4)-(h_2-h_1)}{h_3-h_2}=1-\dfrac{q_L}{q_H}$\\[0.8mm]
$w_{neto}=w_T-w_B,\quad x_4=\dfrac{s_3-s_f}{s_{fg}}$\\[0.6mm]
$h_4=h_f+x_4h_{fg},\quad \dot m=\dot W_{neto}/w_{neto}$
\end{cardidea}
\vspace{1.2mm}
\begin{cardgear}{Ejemplo: 50 MW}{\faCalculator}
$5$ MPa, $400^\circ$C, $P_L=10$ kPa, $\eta_T=0.88$.\\[0.6mm]
$h_3=3196.7$, $h_4=2234.8$, $w_{neto}=955.6$, $q_H=2998.6$ kJ/kg.\\[0.6mm]
$\Rightarrow\ \eta_{th}=31.9\%,\ \dot m=52.3$ kg/s.
\end{cardgear}
\end{column}
\end{columns}
\end{frame}

% ---- Rankine Simple: diagramas ----
\begin{frame}{Rankine Simple --- diagramas de estado}
\begin{columns}[T]
\begin{column}{0.48\textwidth}
\centering
\includegraphics[width=\linewidth,height=3.9cm,keepaspectratio]{rankine_diagrama_Ts}\\[0.5mm]
\fsub{Diagrama $T$-$s$: área encerrada $=w_{neto}$.}
\end{column}
\begin{column}{0.48\textwidth}
\centering
\includegraphics[width=\linewidth,height=3.9cm,keepaspectratio]{rankine_diagrama_Ph}\\[0.5mm]
\fsub{Diagrama $P$-$h$: bomba, caldera, turbina y condensador.}
\end{column}
\end{columns}
\vspace{1mm}
\begin{cardnota}[\faInfoCircle]
\textbf{Calidad a la salida de turbina:} en el ejemplo $x_4=0.854<0.88$ (mínimo ASME). El vapor húmedo erosiona los álabes $\Rightarrow$ motiva el \textbf{recalentamiento}.
\end{cardnota}
\end{frame}

% ---- Rankine con Recalentamiento: definición y etapas ----
\begin{frame}{Rankine con Recalentamiento --- definición y etapas}
\begin{cardidea}{Definición}{\faInfoCircle}
Ciclo de vapor con la expansión \textbf{dividida en dos turbinas} (TAP y TBP) y un \textbf{recalentamiento} intermedio a $P_r$. Permite subir la presión de caldera $P_H$ sin que el vapor salga húmedo, elevando $\eta$ y protegiendo los álabes de la turbina.
\end{cardidea}
\vspace{1.4mm}
\etapastitulo
\begin{etapas}
$1\to2$ & Bombeo isentrópico ($P_L\to P_H$)          & $w_B=h_2-h_1$\\
$2\to3$ & Caldera, isobárica ($P_H$)                  & $q_H=h_3-h_2$\\
$3\to4$ & Expansión en TAP, isentrópica ($\to P_r$)   & $w_{TAP}=h_3-h_4$\\
$4\to5$ & Recalentamiento isobárico ($P_r$)           & $q_{rec}=h_5-h_4$\\
$5\to6$ & Expansión en TBP, isentrópica ($\to P_L$)   & $w_{TBP}=h_5-h_6$\\
$6\to1$ & Condensador, isobárico ($P_L$)              & $q_L=h_6-h_1$\\
\end{etapas}
\end{frame}

% ---- Rankine con Recalentamiento: ecuaciones y ejemplo ----
\begin{frame}{Rankine con Recalentamiento --- ecuaciones y ejemplo}
\begin{columns}[T]
\begin{column}{0.46\textwidth}
\centering
\includegraphics[width=\linewidth,height=4.2cm,keepaspectratio]{rankine_recalentamiento_esquema}\\[0.5mm]
\fsub{Expansión en dos etapas (TAP y TBP) con recalentador.}
\end{column}
\begin{column}{0.5\textwidth}
\begin{cardidea}{Ecuaciones generales}{\faSuperscript}
$w_T=(h_3-h_4)+(h_5-h_6)$\\[0.6mm]
$q_H=(h_3-h_2)+(h_5-h_4)$\\[0.6mm]
$\eta_{th}=\dfrac{w_T-w_B}{q_H},\quad P_r\approx P_H/4$ a $P_H/5$
\end{cardidea}
\vspace{1.2mm}
\begin{cardgear}{Ejemplo: 15 MPa / 500\,$^\circ$C}{\faCalculator}
$P_r=3$ MPa, $T_5=500^\circ$C, $P_L=10$ kPa.\\[0.6mm]
$w_T=1642.0$, $q_H=3727.8$ kJ/kg.\\[0.6mm]
$\Rightarrow\ \eta_{th}=43.6\%,\ x_6=0.878$ (vs.\ $0.760$ sin recalentar).
\end{cardgear}
\end{column}
\end{columns}
\end{frame}

% ---- Rankine Regenerativo: definición y etapas ----
\begin{frame}{Rankine Regenerativo --- definición y etapas}
\begin{cardidea}{Definición}{\faInfoCircle}
Ciclo de vapor que \textbf{extrae una fracción $y$} del vapor de la turbina para precalentar el condensado en un \textbf{calentador abierto} (CAA). Sube la temperatura media de adición de calor y acerca el ciclo al límite de Carnot, reduciendo el calor de entrada.
\end{cardidea}
\vspace{1.4mm}
\etapastitulo
\begin{etapas}
$1\to2$        & Bomba 1 isentrópica ($P_L\to P_e$)        & $w_{B1}=(1-y)(h_2-h_1)$\\
CAA            & Mezcla adiabática (determina $y$)         & $y=\dfrac{h_3-h_2}{h_e-h_2}$\\
$3\to4$        & Bomba 2 isentrópica ($P_e\to P_H$)        & $w_{B2}=h_4-h_3$\\
$4\to5$        & Caldera, isobárica                        & $q_H=h_5-h_4$\\
$5\to e\to6$   & Turbina con extracción ($y$)              & $w_T=(h_5-h_e)+(1-y)(h_e-h_6)$\\
$6\to1$        & Condensador (fracción $1-y$)              & $q_L=(1-y)(h_6-h_1)$\\
\end{etapas}
\end{frame}

% ---- Rankine Regenerativo: ecuaciones y ejemplo ----
\begin{frame}{Rankine Regenerativo --- ecuaciones y ejemplo}
\begin{columns}[T]
\begin{column}{0.46\textwidth}
\centering
\includegraphics[width=\linewidth,height=4.2cm,keepaspectratio]{rankine_regenerativo_esquema}\\[0.5mm]
\fsub{Extracción de fracción $y$ hacia el calentador abierto (CAA).}
\end{column}
\begin{column}{0.5\textwidth}
\begin{cardidea}{Ecuaciones generales}{\faSuperscript}
$y=\dfrac{h_3-h_2}{h_e-h_2}$ (balance del CAA)\\[0.8mm]
$w_T=(h_5-h_e)+(1-y)(h_e-h_6)$\\[0.6mm]
$q_H=h_5-h_4,\quad \eta_{th}=\dfrac{w_T-w_{B1}-w_{B2}}{q_H}$
\end{cardidea}
\vspace{1.2mm}
\begin{cardgear}{Ejemplo: 10 MPa / 500\,$^\circ$C}{\faCalculator}
Extracción a $P_e=1$ MPa, $P_L=10$ kPa.\\[0.6mm]
$y=0.218$, $w_{neto}=1116.1$, $q_H=2600.8$ kJ/kg.\\[0.6mm]
$\Rightarrow\ \eta_{th}=42.9\%$ (vs.\ $40.2\%$ del ciclo simple).
\end{cardgear}
\end{column}
\end{columns}
\end{frame}

% ============================================================
%  SECCIÓN 2 — GAS
% ============================================================
\portadaseccion{2}{Ciclos de Gas --- Brayton}{compresor · cámara · turbina · aire-estándar}{\faIndustry}

% ---- Brayton Simple: definición y etapas ----
\begin{frame}{Brayton Simple --- definición y etapas}
\begin{cardidea}{Definición}{\faInfoCircle}
Ciclo de la \textbf{turbina de gas} (aire-estándar, $c_p$ cte): compresor, cámara de combustión y turbina. Combustión externa (análisis cerrado) o abierta; el parámetro clave es la \textbf{relación de presiones} $r_P=P_H/P_L$. Usado en aviación y generación eléctrica.
\end{cardidea}
\vspace{1.4mm}
\etapastitulo
\begin{etapas}
$1\to2$ & Compresión isentrópica (compresor)        & $w_C=c_p(T_2-T_1),\ T_2=T_1\,r_P^{(\gamma-1)/\gamma}$\\
$2\to3$ & Adición de calor isobárica (cámara)        & $q_H=c_p(T_3-T_2)$\\
$3\to4$ & Expansión isentrópica (turbina)            & $w_T=c_p(T_3-T_4),\ T_4=T_3/r_P^{(\gamma-1)/\gamma}$\\
$4\to1$ & Rechazo de calor isobárico                 & $q_L=c_p(T_4-T_1)$\\
\end{etapas}
\end{frame}

% ---- Brayton Simple: ecuaciones y ejemplo ----
\begin{frame}{Brayton Simple --- ecuaciones y ejemplo}
\begin{columns}[T]
\begin{column}{0.46\textwidth}
\centering
\includegraphics[width=\linewidth,height=4.2cm,keepaspectratio]{brayton_esquema_dispositivos}\\[0.5mm]
\fsub{Compresor, cámara de combustión y turbina en ciclo abierto.}
\end{column}
\begin{column}{0.5\textwidth}
\begin{cardidea}{Ecuaciones generales}{\faSuperscript}
$t=r_P^{(\gamma-1)/\gamma},\ T_2=tT_1,\ T_4=T_3/t$\\[0.6mm]
$w_{neto}=c_p[(T_3-T_4)-(T_2-T_1)]$\\[0.6mm]
$\eta_{th}=1-\dfrac{1}{r_P^{(\gamma-1)/\gamma}},\ \ \mathrm{bwr}=\dfrac{w_C}{w_T}$
\end{cardidea}
\vspace{1.2mm}
\begin{cardgear}{Ejemplo: $r_P=10$}{\faCalculator}
$T_1=300$ K, $T_3=1400$ K, $c_p=1.005$, $\gamma=1.4$.\\[0.6mm]
$T_2=579.2$ K, $T_4=725.1$ K, $w_{neto}=397.7$ kJ/kg.\\[0.6mm]
$\Rightarrow\ \eta_{th}=48.2\%,\ \mathrm{bwr}=41.4\%$.
\end{cardgear}
\end{column}
\end{columns}
\end{frame}

% ---- Brayton Simple: diagrama T-s ----
\begin{frame}{Brayton Simple --- diagrama y ciclo real}
\begin{columns}[T]
\begin{column}{0.47\textwidth}
\centering
\includegraphics[width=\linewidth,height=4.4cm,keepaspectratio]{brayton_diagrama_Ts}\\[0.5mm]
\fsub{Diagrama $T$-$s$: dos isentrópicas y dos isobáricas.}
\end{column}
\begin{column}{0.49\textwidth}
\begin{cardidea}{Ciclo real (eficiencias)}{\faInfoCircle}
$T_{2r}=T_1+\dfrac{T_{2s}-T_1}{\eta_C}$,\ \ $T_{4r}=T_3-\eta_T(T_3-T_{4s})$.\\[0.6mm]
El \textbf{bwr} alto ($\sim40\%$) amplifica las irreversibilidades.
\end{cardidea}
\vspace{1.2mm}
\begin{cardgear}{Ideal vs.\ real ($\eta_C=0.82,\ \eta_T=0.86$)}{\faCalculator}
$w_{neto}$: $397.7\to241.2$ kJ/kg.\\[0.6mm]
$\mathrm{bwr}$: $41.4\%\to58.6\%$.\\[0.6mm]
$\Rightarrow\ \eta_{th}$: $48.2\%\to31.6\%$.
\end{cardgear}
\end{column}
\end{columns}
\end{frame}

% ---- Brayton con Regeneración: definición y etapas ----
\begin{frame}{Brayton con Regeneración --- definición y etapas}
\begin{cardidea}{Definición}{\faInfoCircle}
Brayton al que se añade un \textbf{recuperador} que precalienta el aire comprimido con los gases de escape (requiere $T_4>T_2$). Reduce el combustible necesario para alcanzar $T_3$; ventajoso a \textbf{bajos $r_P$} y alta $T_3/T_1$ (microturbinas, industriales).
\end{cardidea}
\vspace{1.4mm}
\etapastitulo
\begin{etapas}
$1\to2$ & Compresión isentrópica (compresor)         & $w_C=c_p(T_2-T_1)$\\
$2\to5$ & Recuperador, lado frío (isobárico)          & $T_5=T_2+\varepsilon(T_4-T_2)$\\
$5\to3$ & Cámara de combustión (isobárica)            & $q_H=c_p(T_3-T_5)$\\
$3\to4$ & Expansión isentrópica (turbina)             & $w_T=c_p(T_3-T_4)$\\
$4\to6$ & Recuperador, lado caliente (isobárico)      & $T_6=T_4-\varepsilon(T_4-T_2)$\\
$6\to1$ & Rechazo de calor isobárico                  & $q_L=c_p(T_6-T_1)$\\
\end{etapas}
\end{frame}

% ---- Brayton con Regeneración: ecuaciones y ejemplo ----
\begin{frame}{Brayton con Regeneración --- ecuaciones y ejemplo}
\begin{columns}[T]
\begin{column}{0.46\textwidth}
\centering
\includegraphics[width=\linewidth,height=4.2cm,keepaspectratio]{brayton_regeneracion_esquema}\\[0.5mm]
\fsub{Recuperador que precalienta el aire con los gases de escape.}
\end{column}
\begin{column}{0.5\textwidth}
\begin{cardidea}{Ecuaciones generales}{\faSuperscript}
$\varepsilon=\dfrac{T_5-T_2}{T_4-T_2},\quad q_H=c_p(T_3-T_5)$\\[0.8mm]
$\eta_{th,\varepsilon=1}=1-\dfrac{T_1}{T_3}r_P^{(\gamma-1)/\gamma}$ (baja con $r_P$)\\[0.6mm]
Útil si $T_4>T_2$ (a bajos $r_P$).
\end{cardidea}
\vspace{1.2mm}
\begin{cardgear}{Ejemplo: $r_P=6$, $\varepsilon=0.75$}{\faCalculator}
$T_1=300$, $T_3=1200$ K. $T_2=500.4$, $T_4=719.4$, $T_5=664.7$ K.\\[0.6mm]
$q_H=538.0$, $w_{neto}=281.6$ kJ/kg.\\[0.6mm]
$\Rightarrow\ \eta_{th}=52.3\%$ (vs.\ $40.1\%$ del simple).
\end{cardgear}
\end{column}
\end{columns}
\end{frame}

% ---- Comparación ciclos de potencia ----
\begin{frame}{Comparación de ciclos de potencia en $T$-$s$}
\begin{columns}[T]
\begin{column}{0.52\textwidth}
\centering
\includegraphics[width=\linewidth,height=4.7cm,keepaspectratio]{Ts_comparacion_ciclos_potencia}\\[0.5mm]
\fsub{Rankine, Brayton, Otto y Diesel sobre el plano $T$-$s$.}
\end{column}
\begin{column}{0.44\textwidth}
\begin{cardidea}{Idea unificadora}{\faChartLine}
Toda mejora busca \textbf{subir} la temperatura media de adición de calor y \textbf{bajar} la de rechazo, acercándose al rectángulo de Carnot. El área encerrada es $w_{neto}$ y su relación con el calor de entrada define $\eta_t$.
\end{cardidea}
\vspace{1.5mm}
\begin{cardnota}[\faInfoCircle]
\textbf{Recalentamiento} y \textbf{regeneración} recortan las irreversibilidades sin cambiar los focos térmicos.
\end{cardnota}
\end{column}
\end{columns}
\end{frame}

% ============================================================
%  SECCIÓN 3 — COMBUSTIÓN INTERNA
% ============================================================
\portadaseccion{3}{Combustión Interna --- Otto y Diesel}{pistón-cilindro · chispa vs.\ compresión}{\faFire}

% ---- Otto: definición y etapas ----
\begin{frame}{Ciclo Otto --- definición y etapas}
\begin{cardidea}{Definición}{\faInfoCircle}
Ciclo aire-estándar del \textbf{motor de gasolina} (encendido por chispa). La combustión ocurre a \textbf{volumen constante} en el punto muerto superior. Parámetro clave: la relación de compresión $r=V_1/V_2$ ($8$–$11$; más allá aparece detonación).
\end{cardidea}
\vspace{1.4mm}
\etapastitulo
\begin{etapas}
$1\to2$ & Compresión isentrópica              & $T_2=T_1\,r^{\gamma-1},\ P_2=P_1\,r^{\gamma}$\\
$2\to3$ & Adición de calor ($v$ cte, chispa)  & $q_H=c_v(T_3-T_2),\ P_3=P_2\,T_3/T_2$\\
$3\to4$ & Expansión isentrópica (trabajo)     & $T_4=T_3/r^{\gamma-1}$\\
$4\to1$ & Rechazo de calor ($v$ cte, purga)   & $q_L=c_v(T_4-T_1)$\\
\end{etapas}
\end{frame}

% ---- Otto: ecuaciones y ejemplo ----
\begin{frame}{Ciclo Otto --- ecuaciones y ejemplo}
\begin{columns}[T]
\begin{column}{0.44\textwidth}
\centering
\includegraphics[width=\linewidth,height=4.2cm,keepaspectratio]{otto_diagrama_Pv}\\[0.5mm]
\fsub{Diagrama $P$-$v$: combustión isocórica (2$\to$3).}
\end{column}
\begin{column}{0.52\textwidth}
\begin{cardidea}{Ecuaciones generales ($r=V_1/V_2$)}{\faSuperscript}
$q_H=c_v(T_3-T_2),\ q_L=c_v(T_4-T_1)$\\[0.6mm]
$w_{neto}=q_H-q_L,\ T_2/T_1=r^{\gamma-1}$\\[0.6mm]
$\eta_{th}=1-\dfrac{1}{r^{\gamma-1}}$
\end{cardidea}
\vspace{1.2mm}
\begin{cardgear}{Ejemplo: $r=8$}{\faCalculator}
$T_1=300$ K, $q_H=800$ kJ/kg, $\gamma=1.4$.\\[0.6mm]
$T_2=689.2$, $T_3=1803.4$, $T_4=784.9$ K; $P_3=4808$ kPa.\\[0.6mm]
$\Rightarrow\ \eta_{th}=56.5\%,\ w_{neto}=451.8$ kJ/kg.
\end{cardgear}
\end{column}
\end{columns}
\end{frame}

% ---- Diesel: definición y etapas ----
\begin{frame}{Ciclo Diesel --- definición y etapas}
\begin{cardidea}{Definición}{\faInfoCircle}
Ciclo aire-estándar del \textbf{motor de encendido por compresión}. Sin bujía: el aire se comprime hasta la autoignición y el combustible arde a \textbf{presión constante}. Opera a $r$ alto ($14$–$22$) sin detonar; parámetro extra: la relación de corte $r_c=V_3/V_2$.
\end{cardidea}
\vspace{1.4mm}
\etapastitulo
\begin{etapas}
$1\to2$ & Compresión isentrópica                 & $T_2=T_1\,r^{\gamma-1}$\\
$2\to3$ & Adición de calor ($P$ cte, inyección)  & $q_H=c_p(T_3-T_2),\ T_3=T_2\,r_c$\\
$3\to4$ & Expansión isentrópica (trabajo)        & $T_4=T_1\,r_c^{\gamma}$\\
$4\to1$ & Rechazo de calor ($v$ cte)             & $q_L=c_v(T_4-T_1)$\\
\end{etapas}
\end{frame}

% ---- Diesel: ecuaciones y ejemplo ----
\begin{frame}{Ciclo Diesel --- ecuaciones y ejemplo}
\begin{columns}[T]
\begin{column}{0.44\textwidth}
\centering
\includegraphics[width=\linewidth,height=4.2cm,keepaspectratio]{diesel_diagrama_Pv}\\[0.5mm]
\fsub{Diagrama $P$-$v$: combustión isobárica (2$\to$3), corte $r_c$.}
\end{column}
\begin{column}{0.52\textwidth}
\begin{cardidea}{Ecuaciones generales ($r_c=V_3/V_2$)}{\faSuperscript}
$q_H=c_p(T_3-T_2),\ q_L=c_v(T_4-T_1)$\\[0.6mm]
$T_3=T_2r_c,\ T_4=T_1r_c^{\gamma},\ w_{neto}=q_H-q_L$\\[0.6mm]
$\eta_{th}=1-\dfrac{1}{r^{\gamma-1}}\dfrac{r_c^{\gamma}-1}{\gamma(r_c-1)}$
\end{cardidea}
\vspace{1.2mm}
\begin{cardgear}{Ejemplo: $r=18$, $r_c=2$}{\faCalculator}
$T_1=300$ K, $\gamma=1.4$.\\[0.6mm]
$T_2=953.1$, $T_3=1906.2$, $T_4=791.7$ K; $q_H=957.9$ kJ/kg.\\[0.6mm]
$\Rightarrow\ \eta_{th}=63.1\%,\ w_{neto}=604.8$ kJ/kg.
\end{cardgear}
\end{column}
\end{columns}
\end{frame}

% ---- Comparación Otto vs Diesel ----
\begin{frame}{Otto vs.\ Diesel en el plano $P$-$v$}
\begin{columns}[T]
\begin{column}{0.5\textwidth}
\centering
\includegraphics[width=\linewidth,height=4.6cm,keepaspectratio]{CCI_PV_otto_diesel}\\[0.5mm]
\fsub{Combustión isocórica (Otto) vs.\ isobárica (Diesel).}
\end{column}
\begin{column}{0.46\textwidth}
\begin{cardidea}{Clave comparativa}{\faChartLine}
A \textbf{igual $r$}, $\eta_{Otto}>\eta_{Diesel}$ porque $f(r_c)=\frac{r_c^{\gamma}-1}{\gamma(r_c-1)}>1$. Pero el diésel opera a $r$ mucho mayor ($14$--$22$) sin detonar, lo que compensa con creces.
\end{cardidea}
\vspace{1.5mm}
\begin{cardnota}[\faInfoCircle]
Límite $r_c\to1$: el Diesel \textbf{converge al Otto}, $\eta\to1-r^{-(\gamma-1)}$.
\end{cardnota}
\end{column}
\end{columns}
\end{frame}

% ============================================================
%  SECCIÓN 4 — REFRIGERACIÓN
% ============================================================
\portadaseccion{4}{Refrigeración}{compresión de vapor · bomba de calor · COP · R-134a}{\faSnowflake}

% ---- VCR: definición y etapas ----
\begin{frame}{Compresión de Vapor (VCR) --- definición y etapas}
\begin{cardidea}{Definición}{\faInfoCircle}
\textbf{Ciclo Rankine invertido}: consume trabajo para bombear calor de frío a caliente. Fluido: un \textbf{refrigerante} (R-134a). Es la base de refrigeradores, aire acondicionado y chillers. Métrica de desempeño: $\mathrm{COP}_R=q_L/w_C$.
\end{cardidea}
\vspace{1.4mm}
\etapastitulo
\begin{etapas}
$1\to2$ & Compresión isentrópica (compresor)        & $w_C=h_2-h_1$\\
$2\to3$ & Condensador, isobárico (cede $q_H$)        & $q_H=h_2-h_3$\\
$3\to4$ & Válvula de expansión (isoentálpica)        & $h_4=h_3$\\
$4\to1$ & Evaporador, isobárico (absorbe $q_L$)      & $q_L=h_1-h_4=h_1-h_3$\\
\end{etapas}
\end{frame}

% ---- VCR: ecuaciones y ejemplo ----
\begin{frame}{Compresión de Vapor (VCR) --- ecuaciones y ejemplo}
\begin{columns}[T]
\begin{column}{0.46\textwidth}
\centering
\includegraphics[width=\linewidth,height=4.2cm,keepaspectratio]{VCR_esquema_dispositivos}\\[0.5mm]
\fsub{Compresor, condensador, válvula y evaporador.}
\end{column}
\begin{column}{0.5\textwidth}
\begin{cardidea}{Ecuaciones generales}{\faSuperscript}
$\mathrm{COP}_R=\dfrac{q_L}{w_C}=\dfrac{h_1-h_3}{h_2-h_1}$\\[0.8mm]
$x_4=\dfrac{h_3-h_f(P_L)}{h_{fg}(P_L)},\quad q_H=q_L+w_C$\\[0.6mm]
$\mathrm{COP}_{R,C}=\dfrac{T_L}{T_H-T_L}$
\end{cardidea}
\vspace{1.2mm}
\begin{cardgear}{Ejemplo: R-134a, $\dot Q_L=10$ kW}{\faCalculator}
$T_L=-20^\circ$C, $T_H=40^\circ$C. $h_1=386.55$, $h_2=428.35$, $h_3=256.44$ kJ/kg.\\[0.6mm]
$w_C=41.80$, $q_L=130.11$ kJ/kg.\\[0.6mm]
$\Rightarrow\ \mathrm{COP}_R=3.11,\ \dot W_C=3.21$ kW.
\end{cardgear}
\end{column}
\end{columns}
\end{frame}

% ---- VCR: diagrama P-h + flujos ----
\begin{frame}{VCR --- diagrama $P$-$h$ y flujos de energía}
\begin{columns}[T]
\begin{column}{0.48\textwidth}
\centering
\includegraphics[width=\linewidth,height=3.9cm,keepaspectratio]{VCR_diagrama_Ph}\\[0.5mm]
\fsub{Diagrama $P$-$h$: válvula vertical, evaporador horizontal.}
\end{column}
\begin{column}{0.48\textwidth}
\centering
\includegraphics[width=\linewidth,height=3.9cm,keepaspectratio]{refrigeracion_diagrama_flujos}\\[0.5mm]
\fsub{Balance: $q_H=q_L+w_C$ en el ciclo inverso.}
\end{column}
\end{columns}
\vspace{1mm}
\begin{cardnota}[\faInfoCircle]
\textbf{Efecto flash:} el líquido a $P_H$ cruza la válvula a $h$ constante y se vaporiza parcialmente hasta $x_4=\frac{h_3-h_f(P_L)}{h_{fg}(P_L)}=0.389$, enfriándose hasta $T_L$. Límite de Carnot: $\mathrm{COP}_{R,C}=\frac{T_L}{T_H-T_L}=4.22$.
\end{cardnota}
\end{frame}

% ---- Bomba de Calor: definición y etapas ----
\begin{frame}{Bomba de Calor --- definición y etapas}
\begin{cardidea}{Definición}{\faInfoCircle}
\textbf{Mismo ciclo VCR}, pero el efecto útil es el calor $q_H$ \textbf{cedido al espacio caliente} (calefacción); toma calor del exterior, suelo o agua. Un equipo reversible alterna entre refrigerar y calentar. Métrica: $\mathrm{COP}_{HP}=\mathrm{COP}_R+1$ (siempre $>1$).
\end{cardidea}
\vspace{1.4mm}
\etapastitulo
\begin{etapas}
$1\to2$ & Compresión isentrópica (compresor)              & $w_C=h_2-h_1$\\
$2\to3$ & Condensador (cede $q_H$ al interior) \emph{útil} & $q_H=h_2-h_3$\\
$3\to4$ & Válvula de expansión (isoentálpica)             & $h_4=h_3$\\
$4\to1$ & Evaporador (absorbe del exterior)               & $q_L=h_1-h_3$\\
\end{etapas}
\end{frame}

% ---- Bomba de Calor: ecuaciones y ejemplo ----
\begin{frame}{Bomba de Calor --- ecuaciones y ejemplo}
\begin{columns}[T]
\begin{column}{0.44\textwidth}
\centering
\includegraphics[width=\linewidth,height=4.2cm,keepaspectratio]{refrigeracion_diagrama_flujos}\\[0.5mm]
\fsub{Mismo ciclo VCR; el efecto útil es $q_H$ cedido al interior.}
\end{column}
\begin{column}{0.52\textwidth}
\begin{cardidea}{Ecuaciones generales}{\faSuperscript}
$\mathrm{COP}_{HP}=\dfrac{q_H}{w_C}=\dfrac{h_2-h_3}{h_2-h_1}=\mathrm{COP}_R+1$\\[0.8mm]
$\mathrm{COP}_{HP,C}=\dfrac{T_H}{T_H-T_L}$; resistencia: $\mathrm{COP}=1$.
\end{cardidea}
\vspace{0.6mm}
\begin{cardgear}{Ejemplo: R-134a, $\dot W_C=2.5$ kW}{\faCalculator}
{\footnotesize
$T_L'=-5^\circ$C, $T_H'=32^\circ$C. $h_1=396.43$, $h_2=424.86$, $h_3=237.23$ kJ/kg.\\[0.4mm]
$w_C=28.43$, $q_H=187.63$ kJ/kg.\\[0.4mm]
$\Rightarrow\ \mathrm{COP}_{HP}=6.60,\ \dot Q_H=16.5$ kW ($6.6\times$ una resistencia).}
\end{cardgear}
\end{column}
\end{columns}
\end{frame}

% ============================================================
%  CIERRE
% ============================================================
\begin{frame}[plain]
\begin{tikzpicture}[remember picture, overlay]
  \fill[FondoOscuro] (current page.south west) rectangle (current page.north east);
  \node[color=IconoTenue] at ([xshift=-3.4cm,yshift=0.3cm]current page.east)
      {\fontsize{62}{62}\selectfont \faIndustry};
  \fill[GranateVelo] (current page.north west)
      -- ([xshift=8.6cm]current page.north west)
      -- ([xshift=11.3cm]current page.south west)
      -- (current page.south west) -- cycle;
  \fill[GranateOscuro, opacity=0.75] (current page.north west)
      -- ([xshift=6.4cm]current page.north west)
      -- ([xshift=9.1cm]current page.south west)
      -- (current page.south west) -- cycle;
  \draw[Teal, line width=1.5pt] ([xshift=8.6cm]current page.north west)
      -- ([xshift=11.3cm]current page.south west);
  \draw[Naranja, line width=0.9pt] ([xshift=7.4cm]current page.north west)
      -- ([xshift=10.1cm]current page.south west);
  \node[anchor=west, text=white] at ([xshift=1.1cm,yshift=0.75cm]current page.west)
      {\fontsize{24}{28}\selectfont\bfseries ¡Gracias!};
  \draw[white, line width=0.9pt] ([xshift=1.15cm,yshift=0.12cm]current page.west) -- ++(4.8cm,0);
  \node[anchor=west, text=white] at ([xshift=1.15cm,yshift=-0.55cm]current page.west)
      {{\color{white}\faAngleDoubleRight}\hspace{1.5mm}\small Conversión de Energía --- Ciclos \textperiodcentered\ MN121 Termodinámica};
\end{tikzpicture}
\end{frame}

\end{document}
